\section{Actual boundary bookkeeping and empty inputs of the ordered recursion}

\begin{lemma}[Changing the auxiliary ideal never changes the literal geometric boundary]
\label{hir:original-ordered-recursion-literal-boundary}
\uses{hir:original-canonical-ordered-separator-chain,hir:original-canonical-ordered-resolution-tree}
\lean{Hironaka.HironakaOriginalCanonicalEmptyParts.geoTransformAlong_boundary_eq,Hironaka.HironakaOriginalCanonicalEmptyParts.separation_chain_boundary,Hironaka.HironakaOriginalCanonicalEmptyParts.ordered_resolution_boundary}
\leanok
For any actual tower, original marked records with the same initial
ordered boundary have the same literal final boundary. In particular
the recorded final boundary of an actual ordered separation chain or
ordered resolution tree is the boundary obtained by transforming any
original marked record with that same initial list along the same
entire appended tower. No assertion makes the auxiliary ideals equal
or the composite canonical for a different ideal.
\end{lemma}
\begin{proof}
Induct on the literal tower. At length zero the assertion is the given
list equality. At every actual blowup the next boundary is exactly
the current list mapped by the same geometric strict transform, followed
by the inverse image of the same centre. The ideal, carrier and mark
fields do not occur in that formula, so equal current lists give equal
next lists. For a separation chain, induct over its actual terminal
sum stages. Transform composition along the literal append and the
preceding boundary identity identify the next initial list with the
recorded sum-stage boundary; the recursive assertion then gives the
chain's recorded final list. For an ordered resolution tree, apply this
chain result to its actual separator, the same boundary identity to its
actual pure head, and the recursive tree assertion to its actual tail.
Transform composition along both literal appends gives exactly the
claimed whole final boundary of the same original tower.
\end{proof}

\begin{lemma}[An actually empty head makes its entire separation chain empty]
\label{hir:original-empty-head-separation-chain-all-empty}
\uses{hir:original-disjoint-sum-all-empty-canonical-stage}
\lean{Hironaka.HironakaOriginalCanonicalEmptyParts.separation_chain_empty_of_head_empty}
\leanok
If the original head of an actual canonical separation chain is empty,
every actual centre of the entire chain is empty, regardless of its
ordered tail parts. No empty-stage or neutrality certificate is input.
\end{lemma}
\begin{proof}
Induct on the actual chain. Its identity case has no centres. In a
successor, the genuine sum datum contains the vanishing ideal of the
empty original head. The head is disjoint from its current opposite
part, and the genuine canonical empty-sum theorem makes this actual
sum tower all-empty. The strict transform of the actual empty head
is again empty. Apply induction to the unchanged recursive chain.
The literal append of two all-empty histories has every centre empty,
which proves the assertion for the entire same chain.
\end{proof}

\begin{theorem}[An empty filtered input forces the whole actual resolution tree projection-neutral]
\label{hir:original-empty-filtered-parts-tree-projection-neutrality}
\uses{hir:original-ordered-recursion-literal-boundary,hir:original-empty-head-separation-chain-all-empty,hir:original-all-empty-whole-boundary}
\lean{Hironaka.HironakaOriginalCanonicalEmptyParts.ordered_resolution_empty_of_all_parts_empty,Hironaka.HironakaOriginalCanonicalEmptyParts.existsUnique_empty_parts_projection}
\leanok
Let an actual ordered canonical resolution tree have a literally empty
input list after filtering exactly empty carriers. If every original
boundary member is closed, every centre of its same appended tower is
empty, and there is a unique actual biholomorphism to the original base
equal to its projection. Every original closed subset has strict
transform equal to its inverse image. Its recorded whole boundary is
the original inverse-image list followed by precisely one empty entry
per actual blowup; after filtering it is the filtered original list
pulled back under that same biholomorphism. No compactness, ambient
positivity, supplied deletion history or final isomorphism is assumed.
This is the empty-input case, not the entire recursive uniqueness theorem.
\end{theorem}
\begin{proof}
The empty filtered list says every original input carrier is empty.
Induct on the actual resolution tree. Its head separation chain is
all-empty by the preceding lemma. Its head strict carrier remains
empty, so its actual canonical pure-head history has original empty
cosupport and is all-empty as well. Every twice-transformed tail carrier
is still empty; induction makes its actual tail all-empty. The two
literal appends are therefore all-empty. This argument never uses the
permitted bad-set index to infer a centre or a comparison.

An actual blowup of the empty centre is a genuine biholomorphism equal
to its projection. Induction on the just-proved all-empty tower composes
these actual maps, giving its actual biholomorphic projection. Any other
such map has the same underlying projection and hence equals it. A
biholomorphism has no differential exceptional points, and inverse
images of closed sets are closed, so its genuine strict transform is
exactly inverse image for every original closed subset. The first lemma
identifies the tree's recorded final boundary with the boundary of the
same entire tower. The genuine all-empty boundary formula retains each
original member as its inverse image and appends precisely one empty
member per performed blowup. Substitute the actual projection identity
and filter these genuinely empty entries to obtain the stated ordered
filtered identity, without changing the whole unfiltered assertion.
\end{proof}
